\documentclass[11pt]{article}

\usepackage[T1]{fontenc}
\usepackage[utf8]{inputenc}
\usepackage{lmodern}
\usepackage[margin=1.04in]{geometry}
\usepackage{amsmath,amssymb,mathtools,amsthm}
\usepackage{aliascnt}
\usepackage{mathrsfs}
\usepackage{microtype}
\usepackage{enumitem}
\usepackage{xcolor}
\usepackage{booktabs}
\usepackage{url}
\usepackage{hyperref}
\usepackage[capitalize,noabbrev]{cleveref}

\hypersetup{
  colorlinks=true,
  linkcolor=blue!45!black,
  citecolor=green!35!black,
  urlcolor=blue!55!black,
  pdftitle={Exponential Periods and Factorial Moments},
  pdfauthor={Christopher D. Long; Antoine-Auguste Le Blanc},
  pdfsubject={Factorial moments, simplex exponential periods, radial asymptotics, characteristic-five certificates},
  pdfkeywords={factorial conjecture, exponential period, simplex, asymptotic limit, twisted de Rham remainder}
}

\numberwithin{equation}{section}
\setlength{\emergencystretch}{3em}

\newtheorem{theorem}{Theorem}[section]

\newaliascnt{proposition}{theorem}
\newtheorem{proposition}[proposition]{Proposition}
\aliascntresetthe{proposition}

\newaliascnt{lemma}{theorem}
\newtheorem{lemma}[lemma]{Lemma}
\aliascntresetthe{lemma}

\newaliascnt{corollary}{theorem}
\newtheorem{corollary}[corollary]{Corollary}
\aliascntresetthe{corollary}

\theoremstyle{definition}
\newaliascnt{definition}{theorem}
\newtheorem{definition}[definition]{Definition}
\aliascntresetthe{definition}

\newaliascnt{example}{theorem}
\newtheorem{example}[example]{Example}
\aliascntresetthe{example}

\theoremstyle{remark}
\newaliascnt{remark}{theorem}
\newtheorem{remark}[remark]{Remark}
\aliascntresetthe{remark}

\newcommand{\C}{\mathbb C}
\newcommand{\Q}{\mathbb Q}
\newcommand{\Qbar}{\overline{\mathbb Q}}
\newcommand{\R}{\mathbb R}
\newcommand{\Z}{\mathbb Z}
\newcommand{\N}{\mathbb N}
\newcommand{\Aone}{\mathbb A^1}
\newcommand{\Atwo}{\mathbb A^2}
\newcommand{\Gm}{\mathbb G_m}
\newcommand{\Spec}{\operatorname{Spec}}
\newcommand{\Span}{\operatorname{span}}
\newcommand{\trdeg}{\operatorname{trdeg}}
\newcommand{\per}{\operatorname{per}}
\newcommand{\ord}{\operatorname{ord}}
\newcommand{\End}{\operatorname{End}}
\newcommand{\Hom}{\operatorname{Hom}}
\newcommand{\Ext}{\operatorname{Ext}}
\newcommand{\Gal}{\operatorname{Gal}}
\newcommand{\Lie}{\operatorname{Lie}}
\newcommand{\Frac}{\operatorname{Frac}}
\newcommand{\Sym}{\operatorname{Sym}}
\newcommand{\SL}{\operatorname{SL}}
\newcommand{\GL}{\operatorname{GL}}
\newcommand{\Pic}{\operatorname{Pic}}
\newcommand{\Div}{\operatorname{Div}}
\newcommand{\Res}{\operatorname{Res}}
\newcommand{\Tr}{\operatorname{Tr}}
\newcommand{\Crit}{\operatorname{Crit}}
\newcommand{\supp}{\operatorname{supp}}
\newcommand{\rank}{\operatorname{rank}}
\newcommand{\Ecal}{\mathcal E}
\newcommand{\Fcal}{\mathcal F}
\newcommand{\Hcal}{\mathcal H}
\newcommand{\Lcal}{\mathcal L}
\newcommand{\Mcal}{\mathcal M}
\newcommand{\Ocal}{\mathcal O}
\newcommand{\Pcal}{\mathcal P}
\newcommand{\Rcal}{\mathcal R}
\newcommand{\Vcal}{\mathcal V}
\newcommand{\dd}{\,d}
\newcommand{\KZ}{\mathrm{KZ}}

\title{Exponential Periods and Factorial Moments}
\author{Christopher D. Long \and Antoine-Auguste Le Blanc\thanks{Le Blanc cryptographic author identifier: \texttt{b6048125\allowbreak{}30b3535b\allowbreak{}23d6dbb0\allowbreak{}f5abe32d\allowbreak{}1ed1a9de\allowbreak{}35fb4f7d\allowbreak{}04d62cff\allowbreak{}08ea19c8}}}
\date{October 2026}


\newcommand{\Hq}{H_q^2}
\newcommand{\M}{\mathscr M}
\newcommand{\Ncal}{\mathcal N}
\newcommand{\Tcal}{\mathcal T}
\newcommand{\A}{\mathbb A}

\newcommand{\PGL}{\operatorname{PGL}}
\newcommand{\wt}{\operatorname{wt}}
\newcommand{\Bcal}{\mathcal B}

\begin{document}
\maketitle

\begin{center}
\small
\texttt{galizur@gmail.com}
\end{center}

\begin{abstract}
We relate factorial moments of complex polynomials to compact exponential
periods on simplices. An exact radial limit depends only on the highest two
homogeneous layers and tolerates arbitrary complex lower-degree terms.
Using the linear comparison theorem of Paper V, we obtain explicit families
with eventually nonzero factorial moments in every dimension. One family
allows every nonzero algebraic phase parameter, with degree increasing with
dimension. A second family has degree four in every dimension; a
characteristic-five twisted-remainder argument supplies one parameter class
uniformly in the dimension. Exact coprimality calculations remove its
parameter exceptions through dimension 103.

A block-radial limit gives nonflat families with arbitrarily many active
angular variables and arbitrary lower-degree cross-block couplings. Algebraic
specialization extends the Factorial-Conjecture conclusion to the indicated
complex-parameter families; power closure gives infinitely many nonzero
moments there. These conclusions are distinguished from eventual
nonvanishing of moments, which follows from a nonzero limiting period.
No unrestricted Factorial Conjecture or arbitrary-simplex nonvanishing
theorem is asserted.
\end{abstract}

\medskip
\noindent\textit{2020 Mathematics Subject Classification.}
Primary 11J91; Secondary 13N10, 33E30, 41A60.
\tableofcontents

\section{Introduction and scope}
For $n\ge1$ define the factorial functional
\[
 \Lcal_n:\C[x_1,\ldots,x_n]\longrightarrow\C,
 \qquad \Lcal_n(x_1^{e_1}\cdots x_n^{e_n})=\prod_{i=1}^n e_i!.
\]
Equivalently,
\begin{equation}\label{eq:gamma-integral}
 \Lcal_n(f)=\int_{\R_{\ge0}^n} f(x)e^{-(x_1+\cdots+x_n)}\,dx.
\end{equation}
The Factorial Conjecture asks whether
$\Lcal_n(f^m)=0$ for all integers $m\ge1$ forces $f=0$; see
\cite{EdoEssen}. We say that a nonzero polynomial \emph{satisfies the FC
conclusion} when at least one of these moments is nonzero. This paper proves
that conclusion for structured families, not for all polynomials in any new
dimension. It does not address the consecutive-moment assertion of the
Strong Factorial Conjecture.

The principal distinction is between three statements:
\[
 \text{eventually every moment is nonzero}
 \ \Longrightarrow\
 \text{infinitely many are nonzero}
 \ \Longrightarrow\
 \text{the FC conclusion}.
\]
A nonzero radial limit proves the first. Algebraic specialization alone
proves the last; power closure supplies the middle statement in our families.

\paragraph{Results.}
Theorem~\ref{thm:radial} identifies the normalized radial limit.
Theorem~\ref{thm:odd-family} gives an all-dimensional family for every
algebraic scale, using an interior extremum and no critical-root separation.
Theorems~\ref{thm:modfive} and~\ref{thm:elliptic-family} give a quartic family
with a parameter class uniform in dimension. Proposition~\ref{prop:finite}
records the exact finite exception-free range. Theorem~\ref{thm:block-limit}
and Corollary~\ref{cor:block-family} give nonflat analogues. Finally,
Section~\ref{sec:complex} treats complex parameters without claiming
nonzero limiting periods at every transcendental scale.

\paragraph{Dependency.}
Our analytic limits and finite algebraic certificates are proved here.
The arithmetic nonvanishing input is the linear comparison theorem and
elliptic example of Paper V \cite{PaperV}, with the assumptions restated in
Section~\ref{sec:input}. That manuscript is a working research draft;
the present applications do not independently re-prove its geometric and
transcendence arguments. Products of individually nonzero periods below do
not require multicolumn algebraic independence or Paper IV.

Algebraic coefficients belong to $\Qbar$. Write
\[
 S=x_1+\cdots+x_n,\qquad
 \Delta_{n-1}=\{t_i\ge0:\ \textstyle\sum_i t_i=1\}.
\]
The measure $dt=dt_1\cdots dt_{n-1}$ is coordinate measure after eliminating
$t_n$; its volume is $1/(n-1)!$. A zero-dimensional simplex has mass one.
This normalization is used in every slicing and limit formula.

\section{A radial factorial-moment limit}
\begin{theorem}\label{thm:radial}
Let $n,D\ge1$ and
\[
 f=S^D+f_{D-1}+\cdots+f_0,
\]
where $f_j$ is homogeneous of degree $j$, with arbitrary complex
coefficients. Then
\begin{equation}\label{eq:radial-limit}
 \lim_{m\to\infty}\frac{\Lcal_n(f^m)}{\Gamma(Dm+n)}
 =\int_{\Delta_{n-1}}\exp\bigl(f_{D-1}(t)/D\bigr)\,dt.
\end{equation}
If the integral is nonzero, $\Lcal_n(f^m)\ne0$ for every sufficiently large
integer $m$.
\end{theorem}
\begin{proof}
Put $x_i=rt_i$ in \eqref{eq:gamma-integral}. The Jacobian is $r^{n-1}$.
Write $A_j(t)=f_{D-j}(t)$ and expand
\[
 f(rt)=r^D+A_1(t)r^{D-1}+\cdots+A_D(t).
\]
For a multi-index $\mathbf k=(k_1,\ldots,k_D)$ set
$K=\sum_j k_j$ and $J=\sum_j jk_j$. After radial integration the coefficient
of $\prod_j A_j(t)^{k_j}/k_j!$ in the normalized moment is
\[
 C_{m,\mathbf k}=
 \frac{m!}{(m-K)!}\frac{\Gamma(Dm-J+n)}{\Gamma(Dm+n)}
 \quad(K\le m),
\]
and it is zero for $K>m$. Since $K\le J\le DK\le Dm$, every denominator
factor is a positive integer. The first $K$ factors of the gamma ratio's
denominator dominate $m(m-1)\cdots(m-K+1)$, and the remaining factors are
at least one. Thus $0\le C_{m,\mathbf k}\le1$.
For fixed $\mathbf k$,
\[
 C_{m,\mathbf k}\sim D^{-J}m^{K-J}.
\]
Only $k_2=\cdots=k_D=0$ survive, and their sum is $e^{A_1(t)/D}$.
On the compact simplex the absolute terms are bounded by the summable
series $\prod_j\|A_j\|_\infty^{k_j}/k_j!$. Dominated summation, followed by
integration, proves \eqref{eq:radial-limit}. A sequence converging to a
nonzero complex number stays nonzero eventually.
\end{proof}

\begin{remark}\label{rem:scalars}
Adding $D\beta S^{D-1}$, $\beta\in\C$, to the displayed family multiplies
the limit by $e^\beta$. Multiplication of $f$ by $\lambda\ne0$ multiplies
its $m$th moment by $\lambda^m$. Both operations preserve the indicated
nonvanishing conclusions. No sign assumption on the coefficients of $f$
or on its lower layers is used.
\end{remark}

\section{The exponential-period input}\label{sec:input}
For a real algebraic polynomial phase
\[
 q(X,Y)=\alpha X^a+\beta Y^b+q_<,
 \quad\gcd(a,b)=1,\quad a,b\ge2,
\]
assume lower weighted degree for $q_<$ and distinct nondegenerate critical
values. Paper V proves irreducibility of the polynomial twisted module.
For a positive polynomial domain with a strict interior phase minimum or
maximum relative to the specified faces, its linear comparison theorem
implies the following consequence:
\begin{equation}\label{eq:scalar-input}
 0\ne r\in\operatorname{span}_{\Qbar}
 \{X^iY^j:0\le i\le a-2,\ 0\le j\le b-2\}
 \quad\Longrightarrow\quad
 \int_T r e^{\xi q}\,dX\,dY\notin\Qbar
 \quad(\xi\in\Qbar^\times).
\end{equation}
More generally a nonzero \emph{specialized} twisted remainder suffices.
The detector may also be a saddle or the specified continuation of a density
germ. Critical-difference separation and boundary-field disjointness are not
needed for this linear conclusion. Its specialization uses Beukers'
Theorem~3.2 \cite{Beukers}; these stronger hypotheses are used only for
algebraic independence over the boundary field in Paper V.

For the elliptic phase used below, Paper V's continued-density example
provides the same conclusion on the fixed triangle even though its real
saddle is outside that triangle. Only the specified faces are used. The
stronger elliptic independence theorem also implies transcendence over
their value field when the interior remainder is nonzero, but absolute
nonvanishing is sufficient for all moment arguments here.

\section{Every algebraic scale in every dimension}
Fix $n\ge3$ and choose odd $a\ge\max(3,n-1)$. Put
\[
 q_a(X,Y)=X^a-aX+Y^2,
 \qquad T=\{0\le X\le2,\ |Y|\le1-X/2\}.
\]
Its critical points satisfy $X^{a-1}=1$, $Y=0$, and their values
$-(a-1)X$ are distinct. On $T$ the minimum is uniquely attained at $(1,0)$,
strictly below every boundary value. Thus \eqref{eq:scalar-input} applies to
$X^{n-3}$, whose exponent is at most $a-2$.

\begin{theorem}\label{thm:odd-family}
Let $\sigma=x_3+\cdots+x_n$, $D\ge a+1$, $\xi\in\Qbar^\times$, and
$\deg R\le D-2$, with arbitrary complex coefficients in $R$. Define
\begin{align}\label{eq:odd-family}
 f={}&S^D+D\xi\bigl(
 2^a\sigma^a S^{D-1-a}
 -2a\sigma S^{D-2}
 +(x_2-x_1)^2 S^{D-3}\bigr)+R.
\end{align}
Then
\[
 \lim_{m\to\infty}\frac{\Lcal_n(f^m)}{\Gamma(Dm+n)}
 =\frac{1}{2^{n-1}(n-3)!}
   \int_T X^{n-3}e^{\xi q_a(X,Y)}\,dX\,dY\ne0.
\]
Consequently every such $f$ has eventually nonzero factorial moments.
At $\xi=0$ the same conclusion holds with limit $1/(n-1)!$.
\end{theorem}
\begin{proof}
Write $u=t_1$, $v=t_2$. Integrating the remaining $n-2$ simplex coordinates
gives the weight $(1-u-v)^{n-3}/(n-3)!$ on
$u,v\ge0$, $u+v\le1$. The affine map
\[
 X=2(1-u-v),\qquad Y=v-u
\]
has absolute Jacobian four and maps this triangle onto $T$. Thus
\begin{align*}
 &\int_{\Delta_{n-1}}
   e^{\xi q_a(2\sum_{j\ge3}t_j,t_2-t_1)}\,dt\\
 &\hspace{1cm}=
 \frac{1}{2^{n-1}(n-3)!}\int_T X^{n-3}e^{\xi q_a(X,Y)}\,dX\,dY.
\end{align*}
The homogeneous layer of degree $D-1$ in \eqref{eq:odd-family}, divided by
$D$ on $S=1$, is exactly this phase. Apply Theorem~\ref{thm:radial} and
\eqref{eq:scalar-input}. At zero the integrand is one.
\end{proof}

For $a\ge5$ the critical values contain two different opposite pairs,
so ordered critical differences are not separated. This family therefore
uses the weaker linear theorem, not the full-independence theorem disguised
as a scalar corollary. The price of allowing every algebraic scale uniformly
is increasing degree as $n$ increases.

\section{A characteristic-five theorem for the elliptic family}
Take
\[
 q(X,Y)=Y^2-X^3+3X,
 \quad T=\operatorname{conv}\{(-\tfrac12,-1),(\tfrac12,0),(-\tfrac12,1)\},
 \quad \ell=\frac{1-2X-2Y}{4}.
\]
The form $\ell$ is a barycentric coordinate on this triangle. Put
\[
 J_k(z)=\int_T\ell^k e^{zq}\,dX\,dY,
 \qquad M_0=\int_Te^{zq},\qquad M_1=\int_TXe^{zq}.
\]
For $h=1/z$ the twisted operators are
\[
 D_X=h\partial_X+3(1-X^2),\qquad D_Y=h\partial_Y+2Y.
\]
Their additive quotient has basis $1,X$ over $\Z[1/6,h]$. Write
$N_h(r)=a_r(h)+b_r(h)X$ for the unique remainder. The reductions are
\begin{align}
 [X^uY^v]&=[X^{u-2}Y^v]+\frac{h(u-2)}3[X^{u-3}Y^v]
 &&(u\ge2),\label{eq:x-reduction}\\
 [X^uY^v]&=-\frac{h(v-1)}2[X^uY^{v-2}]
 &&(v\ge1),\label{eq:y-reduction}
\end{align}
where zero-coefficient terms are omitted. In particular odd $Y$ powers
reduce to zero. The two reductions commute, and their leading coefficients
are units; the resulting normal form remains free after reduction modulo
five. This is an additive quotient by twisted derivatives, \emph{not} a
quotient algebra.

Set $a_k=a_{\ell^k}$, $b_k=b_{\ell^k}$. Literal Stokes gives
\begin{equation}\label{eq:elliptic-remainder}
 J_k(\xi)=a_k(1/\xi)M_0(\xi)+b_k(1/\xi)M_1(\xi)+B_k(\xi),
\end{equation}
where $B_k(\xi)$ is in the fixed boundary span. A nonzero remainder vector
therefore implies $J_k(\xi)\ne0$ by Paper V.

\begin{theorem}[Uniform nonexactness in characteristic five]\label{thm:modfive}
Let $\xi\in\Qbar^\times$. Suppose that at some nonarchimedean place above five,
\begin{equation}\label{eq:valuation-condition}
 v(\xi)\le0.
\end{equation}
Then $(a_k(1/\xi),b_k(1/\xi))\ne(0,0)$ and $J_k(\xi)\ne0$ for every
integer $k\ge0$.
\end{theorem}
\begin{proof}
In characteristic five, multiplication by $\ell^5$ commutes with both
twisted operators because its two partial derivatives vanish. Unlike
multiplication by $\ell$, it descends to the additive quotient. In the basis
$1,X$, reductions \eqref{eq:x-reduction}--\eqref{eq:y-reduction} give
\[
 A(h)=\begin{pmatrix}h-1&2h^2+2\\2&-h-1\end{pmatrix},
 \qquad \det A(h)=2\quad\text{in }\mathbb F_5[h].
\]
For $v_k=(a_k,b_k)^T$ modulo five the first five vectors are
\begin{align*}
 v_0&=(1,0)^T,&v_1&=(-1,2)^T,&v_2&=(-2h,1)^T,\\
 v_3&=(2h+2,-2h-1)^T,&v_4&=(2h^2+h+1,2h)^T.
\end{align*}
Each is nonzero at every specialization over every extension of
$\mathbb F_5$. For $v_3$ its coordinates sum to one; for $v_4$, the first
plus $(2-h)$ times the second is one modulo five. For $k=5b+r$, $0\le r<5$,
\[
 v_k(h)=A(h)^b v_r(h).
\]
Invertibility proves the assertion for every exponent in characteristic
five. No extrapolation from a finite scan is involved.

Choose a number field containing $\xi$ and the stated valuation. Its element
$h=1/\xi$ is integral at that place, as are the remainder coefficients.
If both specialized coefficients were zero, reduction in the residue field
would contradict the preceding identity. Finally apply
\eqref{eq:elliptic-remainder} and the scalar comparison input.
\end{proof}

Examples include $\xi=i$, $101i$, and $Ni$ for any integer $N$ not divisible
by five. Thus the result includes arbitrarily large purely imaginary phases;
it is not limited to a small-oscillation positivity range.

\section{Degree four in every dimension}
Define the cubic homogeneous polynomial
\begin{equation}\label{eq:Hn}
 H_n=S(x_1+2x_2-S)^2-(x_1-S/2)^3+3S^2(x_1-S/2).
\end{equation}
\begin{theorem}\label{thm:elliptic-family}
Let $n\ge3$, $D\ge4$, and let $\xi\in\Qbar^\times$ satisfy
\eqref{eq:valuation-condition}. For any complex polynomial $R$ of degree at
most $D-2$, put
\begin{equation}\label{eq:elliptic-family}
 f=S^D+D\xi S^{D-4}H_n+R.
\end{equation}
Then
\begin{equation}\label{eq:elliptic-limit}
 \lim_{m\to\infty}\frac{\Lcal_n(f^m)}{\Gamma(Dm+n)}
 =\frac{J_{n-3}(\xi)}{2(n-3)!}\ne0.
\end{equation}
In particular $D=4$ and the same $\xi=101i$ work in every dimension.
\end{theorem}
\begin{proof}
Write $Q(u,v)=q(u-1/2,u+2v-1)$, so $H_n=S^3Q(x_1/S,x_2/S)$.
Simplex slicing gives
\[
 \int_{\Delta_{n-1}}e^{\xi Q(t_1,t_2)}\,dt
 =\frac1{(n-3)!}\int_{u,v\ge0,\,u+v\le1}
 (1-u-v)^{n-3}e^{\xi Q(u,v)}\,du\,dv.
\]
The map $(X,Y)=(u-1/2,u+2v-1)$ has positive Jacobian two, maps the standard
triangle onto $T$, and takes $1-u-v$ to $\ell$. This gives the right side of
\eqref{eq:elliptic-limit}. Theorem~\ref{thm:radial} supplies the limit and
Theorem~\ref{thm:modfive} makes it nonzero.
\end{proof}

For example
\[
 f=S^4+404i\bigl(S(x_1+2x_2-S)^2-(x_1-S/2)^3
                   +3S^2(x_1-S/2)\bigr)+R_{\le2}
\]
has eventually nonzero factorial moments for every $n\ge3$. The quadratic
perturbation may involve every variable with arbitrary complex coefficients.
These are not merely polynomials in three variables with unused coordinates
adjoined.

\section{Finite exceptions and an exact recurrence}\label{sec:exceptions}
For a fixed exponent define
\[
 E_k=\{\xi\in\Qbar^\times:a_k(1/\xi)=b_k(1/\xi)=0\}.
\]
This is finite. At $h=0$ the quotient is reduction modulo $X^2-1,Y$, and
$\ell(1,0)=-1/4$, $\ell(-1,0)=3/4$; hence its remainder vector is not
identically zero for any exponent. Every parameter outside $E_k$ has a
nonzero specialized remainder, whether or not it satisfies the valuation
test. Any member of $E_k$ must be strictly divisible by five at every place
above five. Membership in $E_k$ only defeats this interior-remainder
certificate: it does not prove $J_k(\xi)=0$.

\begin{proposition}[Integer recurrence]\label{prop:recurrence}
Set
\[
 W_k(t)=4^k(a_k(3t),b_k(3t)),\quad W_0=(1,0),\quad W_1=(1,-2).
\]
Its coordinates are in $\Z[t]$, and
\begin{align}\label{eq:recurrence}
 W_{k+2}={}&2W_{k+1}+[3-6(2k+1)t]W_k\\
 &+4ktW_{k-1}-36k(k-1)t^2W_{k-2},\nonumber
\end{align}
with every zero-prefactor term omitted.
\end{proposition}
\begin{proof}
Work coefficientwise in the formal series
$F(s)=N_h(e^{s\ell})$. Twisted integration by parts gives
\[
 N_h(Ye^{s\ell})=\frac{hs}{4}F,\quad
 N_h(X^2e^{s\ell})=(1-hs/6)F,\quad
 N_h(Y^2e^{s\ell})=(h^2s^2/16-h/2)F.
\]
Also $N_h(XYe^{s\ell})=(hs/4)N_h(Xe^{s\ell})$ and
$N_h(Xe^{s\ell})=(1/2-hs/4)F-2F'$. Substituting into $F''$ yields
\[
 F''=(\tfrac12-\tfrac{hs}{4})F'
 +(\tfrac3{16}-\tfrac h8+\tfrac{hs}{48}-\tfrac{h^2s^2}{64})F.
\]
Equate exponential-generating coefficients, put $h=3t$, and multiply by
$4^{k+2}$. This is \eqref{eq:recurrence}; integrality follows by induction.
\end{proof}

\begin{proposition}[Finite exact verification]\label{prop:finite}
For $0\le k\le100$ the two coordinates of $W_k$ have gcd one over $\Q[t]$.
Thus $E_k=\varnothing$ in this range. Consequently
Theorem~\ref{thm:elliptic-family} holds for every nonzero algebraic $\xi$
when $3\le n\le103$. At $\xi=0$ its limit is the simplex volume.
\end{proposition}
\begin{proof}
The supplied script constructs \eqref{eq:recurrence} in exact integer
polynomial arithmetic and performs the 101 Euclidean gcd computations.
All gcd degrees are zero; the maximum coordinate degree is 50. It also
compares the recurrence against independent monomial Stokes reduction
through $k=12$. The recurrence was proved above, so the gcd calculation is
a finite reproducible certificate for exactly the stated range. A constant
gcd excludes simultaneous roots over $\C$, giving the consequence.
\end{proof}

No uniform coprimality assertion for every $k$ is inferred from this test.
The all-exponent theorem is instead the characteristic-five argument of
Theorem~\ref{thm:modfive}.

\section{Block-radial limits and nonflat families}
Partition the variables into nonempty blocks of sizes $n_j$, $1\le j\le r$,
and let $S_j$ denote their sums. Fix integers $d_j\ge1$, and put
\[
 F_0=\prod_{j=1}^r S_j^{d_j},\qquad D=\sum_jd_j,\qquad
 C_d=\prod_jd_j^{d_j}.
\]
A block angular variable $t^{(j)}$ belongs to $\Delta_{n_j-1}$.

\begin{theorem}[Block-radial limit]\label{thm:block-limit}
For a homogeneous $G$ of degree $D-1$ and $\deg R\le D-2$, both with
arbitrary complex coefficients, let $f=F_0+G+R$. Then
\begin{equation}\label{eq:block-limit}
 \lim_{m\to\infty}\frac{\Lcal_n(f^m)}{\prod_j\Gamma(d_jm+n_j)}
 =\int_{\prod_j\Delta_{n_j-1}}
  \exp\!\left(\frac{G(d_1t^{(1)},\ldots,d_rt^{(r)})}{C_d}\right)
  \prod_jdt^{(j)}.
\end{equation}
A nonzero integral again gives eventual nonvanishing.
\end{theorem}
\begin{proof}
Use block polar coordinates $x^{(j)}=s_jt^{(j)}$, whose Jacobian is
$\prod_js_j^{n_j-1}$, and set $s_j=mu_j$. On a fixed compact neighborhood
$V$ of $d=(d_1,\ldots,d_r)$ inside the positive orthant,
\[
 m^{-D}f(mu_1t^{(1)},\ldots,mu_rt^{(r)})
 =\prod_ju_j^{d_j}\left(1+\frac{G(u_1t^{(1)},\ldots,u_rt^{(r)})}
 {m\prod_ju_j^{d_j}}+O(m^{-2})\right),
\]
uniformly in $u\in V$ and the angular variables. Raising the relative
factor to the $m$th power converges uniformly to the corresponding
exponential. Independent gamma random variables of shapes $d_jm+n_j$,
after division by $m$, converge in probability to $d_j$ (their means are
$d_j+n_j/m$ and variances are $d_j/m+n_j/m^2$). Thus the part on $V$ has
the required limit.

We must control the complement without dividing by $F_0$ near a coordinate
boundary. Put $U=\sum_ju_j$ and
\[
 H_m(u,t)=e^{-U}\bigl|m^{-D}f(mu_1t^{(1)},\ldots,mu_rt^{(r)})\bigr|.
\]
On every compact radial set, including its coordinate boundary,
$H_m$ converges uniformly to
$H_0(u)=e^{-U}\prod_ju_j^{d_j}$. This function has a unique global maximum
$M=e^{-D}C_d$ at $d$, and vanishes at the coordinate boundary. Hence on a
compact complement of $V$ its maximum is at most $M e^{-\eta}$ for some
$\eta>0$, for all sufficiently large $m$.

Uniformly for $m\ge1$ there is a constant $C$ with
$H_m\le C e^{-U}(1+U)^D$. For sufficiently large $U$ this is bounded by
$M e^{-\eta}e^{-U/2}$. The resulting tail integral of
$H_m^m\prod_ju_j^{n_j-1}$ is at most a constant times
$M^m e^{-\eta m}m^{-n}$, where $n=\sum_jn_j$.
By Stirling's formula the normalization
$\prod_j\Gamma(d_jm+n_j)/m^{Dm+n}$ is $M^m$ times a nonzero constant and
$m^{-r/2}(1+o(1))$. Both complementary regions are therefore exponentially
negligible after normalization. This proves the limit and justifies the
local gamma argument without an unsafe denominator estimate on the boundary.
\end{proof}

\begin{corollary}[Several active angular blocks]\label{cor:block-family}
For each block choose a polynomial $Q_j$ of degree $a_j\le d_j-1$ and a
scale $\xi_j$ such that
\[
 I_j=\int_{\Delta_{n_j-1}}e^{\xi_jQ_j(t)}\,dt\ne0.
\]
Let $\widetilde Q_j(x^{(j)})=S_j^{a_j}Q_j(x^{(j)}/S_j)$, and set
\[
 G=\sum_j d_j\xi_j\left(\prod_{l\ne j}S_l^{d_l}\right)
 S_j^{d_j-1-a_j}\widetilde Q_j(x^{(j)}).
\]
For every $\deg R\le D-2$, the polynomial $f=F_0+G+R$ has limiting ratio
$\prod_jI_j\ne0$ in \eqref{eq:block-limit}. In particular the odd-degree and
elliptic constructions may be assigned independently to arbitrarily many
blocks. Inactive singleton blocks contribute mass one.
\end{corollary}
\begin{proof}
At $x^{(j)}=d_jt^{(j)}$, the quotient $G/C_d$ is
$\sum_j\xi_jQ_j(t^{(j)})$. The integral in \eqref{eq:block-limit} factors.
\end{proof}

The top $F_0$ is genuinely nonflat when there are at least two blocks.
Lower-degree cross-block couplings in $R$ are unrestricted. The degree-$D-1$
correction in this corollary is deliberately block-separable; arbitrary
coupled angular phases at that degree would require a new nonvanishing
argument. No several-phase comparison theorem is used to multiply the
individual nonzero numbers.

\section{Complex parameters and power closure}\label{sec:complex}
\begin{lemma}[Algebraic specialization of moment equations]\label{lem:specialization}
Let a polynomial family $f_\theta$ be defined over $\Qbar$ on an affine
finite-type parameter space, and let $U$ be a Zariski-open subset defined
over $\Qbar$. If no $\Qbar$-point of $U$ has all its positive factorial moments zero,
then no $\C$-point of $U$ does.
\end{lemma}
\begin{proof}
A hypothetical complex point lies in a principal affine open contained in
$U$. Adjoin an inverse to its defining nonvanishing polynomial. The moment
polynomials generate an ideal in this finitely generated $\Qbar$-algebra.
The complex point makes it proper. The weak Nullstellensatz gives a
$\Qbar$-valued point of its zero set, at which every moment still vanishes and
the same open condition holds, a contradiction. No finite uniform moment
bound is required: Noetherianity applies to the ideal of all moments.
\end{proof}

\begin{theorem}\label{thm:complex}
The odd-degree family \eqref{eq:odd-family} satisfies the FC conclusion for
every $\xi\in\C$. The elliptic family \eqref{eq:elliptic-family} satisfies it
for every $\xi\in\C\setminus E_{n-3}$, with zero included. For
$3\le n\le103$ it therefore holds for every complex scale. In all these
families there are infinitely many nonzero factorial moments.
Analogous conclusions hold for the block-separable families when each scale
lies outside its indicated finite algebraic exception set.
\end{theorem}
\begin{proof}
For the odd-degree family all algebraic scales give a nonzero limit by
Theorem~\ref{thm:odd-family}, including zero. Apply
Lemma~\ref{lem:specialization}, allowing the coefficients of $R$ as additional
parameters. For the elliptic family at a nonzero scale outside $E_{n-3}$,
at least one specialized remainder coordinate is nonzero. After clearing
powers of $\xi$, invert that coordinate and $\xi$. At every algebraic point
in this open set, scalar comparison and Theorem~\ref{thm:radial} give a
nonzero limit. The lemma applies. At $\xi=0$ the radial limit is directly
the simplex volume, for arbitrary complex $R$. The block argument uses the
product of the selected nonvanishing conditions and Corollary~\ref{cor:block-family}.

These families are closed under positive powers with the same angular phase
and scales. For a flat top,
\[
 (S^D+f_{D-1}+R)^N
 =S^{DN}+N S^{D(N-1)}f_{D-1}+R',\qquad \deg R'\le DN-2.
\]
Thus $D$ changes to $DN$ while $f_{D-1}/D$ on $S=1$ is unchanged.
For the block family $d_j$ changes to $Nd_j$ and $G$ to
$NF_0^{N-1}G$, preserving the same block angular correction.
If only finitely many moments of $f$ were nonzero, choose $N$ larger than
all their indices. Then every positive moment of $f^N$ would vanish,
contradicting its FC conclusion.
\end{proof}

\begin{remark}
The specialization theorem does not show that the limiting exponential
period is nonzero for every transcendental scale. Its conclusion concerns
the polynomial moment equations. Eventual nonvanishing is asserted where a
nonzero limit was proved; infinitely many nonzero moments is the conclusion
obtained from specialization and power closure alone.
\end{remark}

\section{General transfers, verification, and remaining problems}
\begin{proposition}[A finite exceptional set for each slicing exponent]
\label{prop:transfer}
Suppose a triangular phase satisfies Paper V's scalar comparison and has
a Laurent Stokes basis with $h=1/z$ reduction regular at $h=0$, specializing
to its Jacobian quotient. Let $\ell$ be a barycentric linear form which is
nonzero at at least one critical point. For each $k\ge0$, the integral
$\int_T\ell^k e^{\xi q}$ is nonzero for all but finitely many
$\xi\in\Qbar^\times$.
\end{proposition}
\begin{proof}
The Jacobian class of $\ell^k$ is nonzero, by evaluation at that critical
point. Thus its polynomial remainder vector is not identically zero in
$h$. Outside the finite common-zero set its specialization is nonzero;
apply scalar comparison. A hypothetical exceptional parameter would only
be a failure of this certificate, not a proven zero of the integral.
\end{proof}

The characteristic-five theorem adds a uniform arithmetic certificate for
all exponents to this finite-exception argument. More generally, in good
characteristic $p$, multiplication by a $p$th power commutes with the
twisted derivatives. An invertible induced matrix and nonzero initial
vectors for exponents $0,\ldots,p-1$ give an all-exponent nonexactness test.
It must be proved in the actual free twisted quotient; multiplication by an
arbitrary amplitude need not descend.

\paragraph{Verification.}
Run the exact script from the repository root:
\begin{verbatim}
python3 scripts/check_comparison_factorial.py \
    --max-k 100 --output .build/comparison-factorial.json
\end{verbatim}
It checks the two-coordinate mod-five matrix, its unit determinant and five
initial vectors, the integer recurrence against direct normal forms, all 101
finite gcds, affine Jacobians and low-degree normalization identities, and
related Paper V examples. The tests use exact arithmetic, not oscillatory
quadrature. The finite range in Proposition~\ref{prop:finite} is
computational; Theorem~\ref{thm:modfive} is an all-exponent proof.

\paragraph{Boundaries.}
The explicit families do not establish unrestricted FC$(n)$, arbitrary
flat-top phases, or unrestricted simplex exponential nonvanishing. Neither
does a zero leading radial limit disprove FC: higher asymptotic coefficients
may be nonzero. Uniform emptiness of every $E_k$ and nonvanishing for coupled
angular phases remain separate targets. The companion simplex candidate
\cite{SimplexCandidate} records these distinctions and a counterexample to
the overly broad positive-amplitude induction. The multicolumn project is
separate from the product argument of Corollary~\ref{cor:block-family}.

\paragraph{AI assistance and responsibility.}
OpenAI ChatGPT and Anthropic Claude were used for proof exploration,
drafting, symbolic checks, and adversarial review. They are not authors.
The authors are responsible for the statements, proofs, citations, and
remaining errors. Compilation and finite certificates do not constitute
independent peer review or formal verification.

\small
\begin{thebibliography}{99}
\bibitem{Beukers} F.~Beukers,
\emph{A refined version of the Siegel--Shidlovskii theorem},
Ann. of Math. (2) \textbf{163} (2006), 369--379.
\href{https://doi.org/10.4007/annals.2006.163.369}{doi:10.4007/annals.2006.163.369}.

\bibitem{EdoEssen} E.~Edo and A.~van den Essen,
\emph{The Strong Factorial Conjecture},
preprint, 2013, \href{https://arxiv.org/abs/1304.3956}{arXiv:1304.3956}.
We use its definition of the Factorial Conjecture, not a Strong-Factorial
conclusion.

\bibitem{PaperV} C.~D.~Long and A.-A.~Le~Blanc,
\emph{Relative Exponential Periods on Polynomial Chains:
Algebraic Independence over Fixed Boundary Fields}, Paper~V,
September 2026, two-level revision October 2026, working manuscript.
\href{https://github.com/long-mathematics/exponential-integral-theorem/blob/main/papers/05-polynomial-flags/polynomial-flag-exponential-periods.tex}{Repository source}.

\bibitem{SimplexCandidate} C.~D.~Long and A.-A.~Le~Blanc,
\emph{Simplex Exponential Periods: Reductions and Open Comparison Problems},
unnumbered candidate note, October 2026, in the same repository.
\end{thebibliography}
\end{document}
